% Options for packages loaded elsewhere
% Options for packages loaded elsewhere
\PassOptionsToPackage{unicode}{hyperref}
\PassOptionsToPackage{hyphens}{url}
\PassOptionsToPackage{dvipsnames,svgnames,x11names}{xcolor}
%
\documentclass[
  letterpaper,
  DIV=11,
  numbers=noendperiod,
  oneside]{scrartcl}
\usepackage{xcolor}
\usepackage[left=1in,marginparwidth=2.0666666666667in,textwidth=4.1333333333333in,marginparsep=0.3in]{geometry}
\usepackage{amsmath,amssymb}
\setcounter{secnumdepth}{-\maxdimen} % remove section numbering
\usepackage{iftex}
\ifPDFTeX
  \usepackage[T1]{fontenc}
  \usepackage[utf8]{inputenc}
  \usepackage{textcomp} % provide euro and other symbols
\else % if luatex or xetex
  \usepackage{unicode-math} % this also loads fontspec
  \defaultfontfeatures{Scale=MatchLowercase}
  \defaultfontfeatures[\rmfamily]{Ligatures=TeX,Scale=1}
\fi
\usepackage{lmodern}
\ifPDFTeX\else
  % xetex/luatex font selection
\fi
% Use upquote if available, for straight quotes in verbatim environments
\IfFileExists{upquote.sty}{\usepackage{upquote}}{}
\IfFileExists{microtype.sty}{% use microtype if available
  \usepackage[]{microtype}
  \UseMicrotypeSet[protrusion]{basicmath} % disable protrusion for tt fonts
}{}
\makeatletter
\@ifundefined{KOMAClassName}{% if non-KOMA class
  \IfFileExists{parskip.sty}{%
    \usepackage{parskip}
  }{% else
    \setlength{\parindent}{0pt}
    \setlength{\parskip}{6pt plus 2pt minus 1pt}}
}{% if KOMA class
  \KOMAoptions{parskip=half}}
\makeatother
% Make \paragraph and \subparagraph free-standing
\makeatletter
\ifx\paragraph\undefined\else
  \let\oldparagraph\paragraph
  \renewcommand{\paragraph}{
    \@ifstar
      \xxxParagraphStar
      \xxxParagraphNoStar
  }
  \newcommand{\xxxParagraphStar}[1]{\oldparagraph*{#1}\mbox{}}
  \newcommand{\xxxParagraphNoStar}[1]{\oldparagraph{#1}\mbox{}}
\fi
\ifx\subparagraph\undefined\else
  \let\oldsubparagraph\subparagraph
  \renewcommand{\subparagraph}{
    \@ifstar
      \xxxSubParagraphStar
      \xxxSubParagraphNoStar
  }
  \newcommand{\xxxSubParagraphStar}[1]{\oldsubparagraph*{#1}\mbox{}}
  \newcommand{\xxxSubParagraphNoStar}[1]{\oldsubparagraph{#1}\mbox{}}
\fi
\makeatother


\usepackage{longtable,booktabs,array}
\newcounter{none} % for unnumbered tables
\usepackage{calc} % for calculating minipage widths
% Correct order of tables after \paragraph or \subparagraph
\usepackage{etoolbox}
\makeatletter
\patchcmd\longtable{\par}{\if@noskipsec\mbox{}\fi\par}{}{}
\makeatother
% Allow footnotes in longtable head/foot
\IfFileExists{footnotehyper.sty}{\usepackage{footnotehyper}}{\usepackage{footnote}}
\makesavenoteenv{longtable}
\usepackage{graphicx}
\makeatletter
\newsavebox\pandoc@box
\newcommand*\pandocbounded[1]{% scales image to fit in text height/width
  \sbox\pandoc@box{#1}%
  \Gscale@div\@tempa{\textheight}{\dimexpr\ht\pandoc@box+\dp\pandoc@box\relax}%
  \Gscale@div\@tempb{\linewidth}{\wd\pandoc@box}%
  \ifdim\@tempb\p@<\@tempa\p@\let\@tempa\@tempb\fi% select the smaller of both
  \ifdim\@tempa\p@<\p@\scalebox{\@tempa}{\usebox\pandoc@box}%
  \else\usebox{\pandoc@box}%
  \fi%
}
% Set default figure placement to htbp
\def\fps@figure{htbp}
\makeatother





\setlength{\emergencystretch}{3em} % prevent overfull lines

\providecommand{\tightlist}{%
  \setlength{\itemsep}{0pt}\setlength{\parskip}{0pt}}



 


\KOMAoption{captions}{tableheading}
\makeatletter
\@ifpackageloaded{caption}{}{\usepackage{caption}}
\AtBeginDocument{%
\ifdefined\contentsname
  \renewcommand*\contentsname{Table of contents}
\else
  \newcommand\contentsname{Table of contents}
\fi
\ifdefined\listfigurename
  \renewcommand*\listfigurename{List of Figures}
\else
  \newcommand\listfigurename{List of Figures}
\fi
\ifdefined\listtablename
  \renewcommand*\listtablename{List of Tables}
\else
  \newcommand\listtablename{List of Tables}
\fi
\ifdefined\figurename
  \renewcommand*\figurename{Figure}
\else
  \newcommand\figurename{Figure}
\fi
\ifdefined\tablename
  \renewcommand*\tablename{Table}
\else
  \newcommand\tablename{Table}
\fi
}
\@ifpackageloaded{float}{}{\usepackage{float}}
\floatstyle{ruled}
\@ifundefined{c@chapter}{\newfloat{codelisting}{h}{lop}}{\newfloat{codelisting}{h}{lop}[chapter]}
\floatname{codelisting}{Listing}
\newcommand*\listoflistings{\listof{codelisting}{List of Listings}}
\makeatother
\makeatletter
\makeatother
\makeatletter
\@ifpackageloaded{caption}{}{\usepackage{caption}}
\@ifpackageloaded{subcaption}{}{\usepackage{subcaption}}
\makeatother
\makeatletter
\@ifpackageloaded{sidenotes}{}{\usepackage{sidenotes}}
\@ifpackageloaded{marginnote}{}{\usepackage{marginnote}}
\makeatother
\usepackage{bookmark}
\IfFileExists{xurl.sty}{\usepackage{xurl}}{} % add URL line breaks if available
\urlstyle{same}
\hypersetup{
  pdftitle={5: Socialization and Interaction},
  colorlinks=true,
  linkcolor={blue},
  filecolor={Maroon},
  citecolor={Blue},
  urlcolor={Blue},
  pdfcreator={LaTeX via pandoc}}


\title{5: Socialization and Interaction}
\author{}
\date{}
\begin{document}
\maketitle


\section{Socialization}\label{socialization}

In the \href{../chapters/Ch_4_Meaning.qmd}{last unit}, we established
how meaning is one of the central components of sociology, and
particularly its components, of values, norms, and material culture. It
remains to be seen, however, how people come to know these elements of
society. This process, in turn, is called \textbf{socialization}.

\marginnote{\begin{footnotesize}

\textbf{Socialization} is the ongoing process, beginning at birth,
through which people develop an awareness of social norms and values and
achieve a distinct sense of self.

\end{footnotesize}}

One key out put of the process of socialization is the distinctive
senses of self that we (and, at a more macro-level, groups) develop
relative to society at large and other people and groups. For instance
(and as we will discuss later in the course), to answer the question
``what makes you \emph{you},'' in contemporary developed societies we
would typically answer with a variety of markers---our gender, our race,
perhaps our religion, our national origin, and so on. We call these
markers of who we are \textbf{identity}, and we furtner recognize that
the concept contains a dynamic between two elements: what people think
is categorically distinctive about us (\textbf{social identities}
ascribed to us) and how those categories intersect, or stand in tension
with, our senses of ``who we really are'' (our \textbf{self-identity}).

\marginnote{\begin{footnotesize}

\textbf{Identity} is the distinctive intersection for a person or group
of who they are and what is meaningful to them. It is composed of the
dynamic between \textbf{social-identities}, the significant categories
attributed to people and groups by others, and \textbf{self-identity},
which is how those categories affect a person or groups' understanding
of themselves.

\end{footnotesize}}

As if that were not complex enough, our identities and socialization
often relate, in turn, to our \textbf{social roles.} As noted by Erving
Goffman (about whom more later in this unit), a social role has a close
connection with the idea of a theatrical role---there are expectations
about how you should act and react, in a given context, if you are to
persuasively ``give off'' the markers of a good performance of that
role. Think of being a part of your family (another key social- and
self-identity for many people): there are expectations for how parents
and children should behave, and those are related to one another; a
parent \emph{should} guide a child's development, and a child
\emph{should} explore the world in light of that guidance but also show
devotion and perhaps support to their parent.

\marginnote{\begin{footnotesize}

\textbf{Social roles} are defined expectations for behavior and social
identity for a person given by a person's social position and the social
setting for an interaction.

\end{footnotesize}}

In the context of our class together (the relevant \textbf{social
setting}), I have the social role of a professor, which means that you
can expect me to provide you a coherent introduction to the material
announced in the course, and I can expect you as students (another
social role) to try to understand the material, ask good-faith relevant
questions, and so on. Likewise, I can depend on my Teaching Assistants
to perform a role in the setting.

\marginnote{\begin{footnotesize}

The \textbf{social setting} describes the time, place, and role
expectations for people within an interaction.

\end{footnotesize}}

\subsection{Why Is Socialization
Necessary?}\label{why-is-socialization-necessary}

There is some speculation that socialization has biological roots,
specifically the
\href{https://en.wikipedia.org/wiki/Obstetrical_dilemma}{obstetrical
dilemma}, which limits the maximum brain size at birth for bipedal
primates.\sidenote{\footnotesize This hypothesis has been questioned, but it remains
  true that regardless of the reasons for it, humans experience
  \href{https://pmc.ncbi.nlm.nih.gov/articles/PMC6870052/}{rapid
  increases in brain volume}, which are much higher than primates and
  other mammals.} By consequence of this fact, and also the sheer
technical and cultural complexity of social life, it is simply
impossible for people, at birth, to have all relevant cultural
information necessary to exist as they do stored in their brains at the
time of birth.

\subsection{Who Socializes?}\label{who-socializes}

There are numerous groups who socialize you. Above all, there are those
who provide your \textbf{primary socialization} in early childhood until
your adolescence. These usually include your family and (it's in the
name!) primary caregiver (whether a biological parent or otherwise) who
teach you the initial stages of how to ``be'' in the world. This often
includes teaching you to talk, care for yourself in a basic
way\sidenote{\footnotesize Yes, this includes potty training!}, and restrain basic
urges in the name of social norms.

To get a sense of what primary socialization looks like in practice,
think of dinner at a ``family'' restaurant. If you've gone with
\emph{your} family to such a restaurant, you know that it's not
\emph{always} an experience that's so joyous you lose track of time.
Maybe you don't like your cousin or sibling, or maybe a grandparent told
you a story for the 30th time or commented about your appearance. You
may have \emph{wanted} to get up and storm out, pull out your phone, or
even scream at the person, but you (hopefully) didn't, because you were
``acting like an adult.''

Compare this to the small children in other families at the
restaurant.\sidenote{\footnotesize Or maybe your family has small children in it. I
  don't know your life.} In those cases, the children---because their
primary socialization is not yet complete---tend to behave \emph{very}
differently. They actually act on those impulses, and may even throw
food, break down crying, or get out of their seats and start crawling
around or running between tables. And, assuming that the \emph{parents}
are normally socialized, you will witness primary socialization
happening. ``Timmy, shush while your nonna is talking,'' a parent might
say to a child, ``We don't interrupt!'' Or, ``it's not polite to spit
out food you don't like!'' Or, ``If you keep crying, we'll leave.'' Or,
``Make sure that you put your napkin on your lap, so that food doesn't
get on your clothes!'' All of those are examples of when social norms,
values, and information about material culture is being directly and
explicitly transmitted to children.

Running congruently with primary socialization as children grow into
adulthood, \textbf{secondary socialization} is undertaken by media,
family, work, and peer-groups. Vehicles of secondary socialization
typically augment those of primary socialization, but also include those
reaching far beyond an individual's direct sphere of ordinary
interaction and including forms of socialization that a person might
directly seek out for themselves.

\subsubsection{Facing Primary
Socialization}\label{facing-primary-socialization}

Here are two excellent examples of primary socialization.

The first is the famous
\href{https://www.youtube.com/watch?v=f1Jw0-LExyc}{still face
experiment}. Here, we can see the sheer intensity of socializing and
cueing by young children. The apparatus of the experiment is for a
mother to simply cease all interaction and hold a ``still face'' in
front of their child. This is often taught as an example of interaction,
but it's an equally powerful example of \emph{socialization}, since the
act of interaction is the reward and punishment system of socialization.
For instance, in the absence of positive reward from the mother (in the
form of smiles, attention, and so on), one can easily observe the
child's behavior literally regressing and becoming \emph{less}
socialized.

The second is a remarkable
\href{https://www.tiktok.com/@brandz100/video/7345831569459825950}{example}
of primary socialization in early-childhood education. In this case, the
educator uses deft attention and explicit rewards (in the form of
high-fives) to explicitly socialize young children (my guess is around
4) into collectively listening to him. Here, you can \emph{also} see how
social coordination is a vehicle for collective socialization; the
children are not only experiencing bilateral interaction (between each
of them and the teacher), they are also being socialized by the behavior
of their peers (the other children). Further, it is worth noting that
the inciting incident for the interaction that was recorded was that
other children were going outside to play, and the children crying at
the start of the video were left behind because they were misbehaving.

\subsubsection{Ongoing Socialization}\label{ongoing-socialization}

As noted above, secondary socialization continues throughout the life
course and is differentiated from primary socialization because it
typically includes a much wider set of social influences than primary
socialization and also includes forms of socialization that people seek
out for themselves.

Examples of this ongoing socialization abound. Consider just a few:

\begin{itemize}
\tightlist
\item
  a joke-y set of explanations about
  \href{https://www.youtube.com/watch?v=aQLOmKkEKYM}{how to greet
  co-workers depending on familiarity}
\item
  a website that explains
  \href{https://www.stonybrook.edu/mobility-and-parking/parking/students.html}{how
  and where to park on campus}
\item
  the
  \href{https://www.investor.gov/introduction-investing/investing-basics/investment-products/stocks}{federal
  government explaining what stocks are and how to buy and sell them}
\item
  \href{https://www.reddit.com/r/islam/comments/14z0u77/whats_the_etiquette_for_being_in_a_mosque/}{rules
  of etiquette for how to behave in a mosque}.\sidenote{\footnotesize I could multiply
    these examples literally a thousand-fold in a very short time. I
    chose these because they show some of the micro- macro-sociological
    scope of ongoing socialization.}
\end{itemize}

\subsection{The Life Course: Categories of
Socialization}\label{the-life-course-categories-of-socialization}

As the depth and breadth of ongoing socialization implies, there are a
vast web of social settings and subcultures that one may need
socialization into. However, it is \emph{also} the case that these
settings and subcultures tend to be grouped together in terms of where
we stand in our \textbf{life course}, which is to say, the
socially-meaningful stages of our lives which tend to correlate with
typical interlaced social settings, role expectations, and identities.

Think of moving from high-school to college. Depending on your specific
circumstances, this could have involved

\begin{itemize}
\tightlist
\item
  Moving from home, sometimes a fairly long distance, to come to Stony
  Brook.
\item
  Ending or starting a romantic relationship.
\item
  Learning a variety of new skills (for instance, your own laundry) and
  to live on your own, or with a roommate.
\item
  Adjusting to a new pace and style of learning.
\end{itemize}

We call the connection among these various new social settings, and the
socialization that goes with them, a life-course transition.

Such life-course transitions are not limited to your early life;
instead, the extend throughout you life. Indeed, it is a common
situation for parents (to flip the above example about going to college
on its head) to experience ``empty-nesting,'' in which \emph{they} must
learn how to be older adults when their children are out of the house
and no longer need such intensive parenting.

Finally, stages in the life course (what you transition into and out of)
are themselves social constructions. That means that the stages of
``childhood,'' ``adulthood,'' and ``retirement/old age,'' for instance,
are not natural or biologically fore-ordained, and instead have varied
quite considerably. To get a sense of this, consider that when the
United States implemented old-age pensions (called Social Security, and
intended to cut poverty among the retired) and the first monthly Social
Security Check
\href{https://en.wikipedia.org/wiki/Social_Security_(United_States)}{was
issued in 1940}, average life expectancy was actually \emph{lower} than
the retirement age (average life expectancy was 61.4 for men and 65.7
for women and the retirement age in 1940 was 65.5).\sidenote{\footnotesize Source:
  \href{https://www.ssa.gov/oact/TR/TR01/lr5A3-h.html}{Social Security
  Administration}.} By contrast, today, the retirment age
\href{https://www.ssa.gov/oact/progdata/nra.html}{is 67}, and average
life expectancy in the US in 2024 was
\href{https://www.cdc.gov/nchs/products/databriefs/db548.htm}{79}. This
means, simply, that there is a group of people who now live on their
savings and national pensions for about a decade, on average, whom we
call ``retirees,'' which didn't exist with the same structure in your
great-grandparents' time.\sidenote{\footnotesize A little more precisely, we want to
  know the average life expectancy \emph{at the age of retirement}, not
  at birth. Today, at retirement age, that life expectancy is about 15
  years for men and about 18 years for women.
  (\href{https://www.health.ny.gov/health_care/medicaid/publications/docs/adm/06adm-5att8.pdf}{Source}.)}

\subsection{Theories of Socialization}\label{theories-of-socialization}

There are two general theories of socialization relevant for many
sociologists, and one rich, applied example of their combination, in the
work of Howard Becker.

\subsubsection{Mead}\label{mead}

First, the sociologist
\href{https://en.wikipedia.org/wiki/George_Herbert_Mead}{George Herbert
Mead}, building from the classical sociology of (among others) Emile
Durkheim, proposed that socialization was a matter of the tension
between people's individual impulses (which Mead called the ``I'')
regardless of social context, and their senses of themselves given the
social identities attributed to them by others and the reactions by
others to one's own behavior (called the ``me''). For Mead, the process
of socialization was the progressive incorporation of the ``me'' into
peoples' identities, and, therefore, their incorporation of the view of
others into the view of themselves.

For Mead, this process can be directly observed in things like
children's games. Playing with a very young child, it is common for the
child to rapidly (and selfishly) change the terms of the game, or even
its substance, over the course of play. However, as children get older
and socialization progresses, games become more robust and durable, to
the point where they overtake any individual child's desires. Put
differently, a point comes where the game continues whether or not any
individual playing it would like to be playing something else, or
nothing at all.

Finally, Mead also noted that the process of socialization was one
whereby the concrete ``others'' referenced in the construction of the
``me'' became increasingly abstract over the course of socialization.
During primary socialization, others are usually concrete people---thus,
for instance, we wonder ``what would my mother think of me doing this?''
as our ``me'' becomes more salient to us. As we pass to secondary
socialization, meanwhile, the others become a general picture of what
``people'' would think of our behaviors, becoming, in Mead's terms,
``the generalized other.''

\subsubsection{Piaget}\label{piaget}

The second major theorist of socialization is
\href{https://en.wikipedia.org/wiki/Jean_Piaget}{Jean Piaget}, who, like
Mead, emphasized that socialization is the process of increasingly
abstract reasoning. But \emph{unlike} Mead, Piaget's theory only
indrectly gestured to socialization by other people, instead stressing
direct interaction with the world.

For Piaget, this meant that children pass through several phases, namely

\begin{itemize}
\tightlist
\item
  The sensorimotor phase, in which children learn the boundaries between
  themselves and the rest of the world (primarily by touching);
\item
  Preoperational, in which children make disorderdered symbolic
  representation (babbling before speaking coherently) and egocentric
  speach (confusing ``I'' and ``you,'' for instance);
\item
  Concrete operational, in which children gain the ability for basic
  symbolic interaction and initial abstraction, such as being able to
  represent things with numbers and manipulate those numbers without
  directly interacting with concrete reprentations;\sidenote{\footnotesize Concretely,
    this means being able to answer a question like ``What is ten minus
    five'' without counting the numbers on the digits of your hand.}
\item
  Formal operational, in which people become capable of highly-abstract
  forms of reasoning and intellectual manipulation of higher-order
  intellectual problems that cannot be materially represented.
\end{itemize}

\subsection{\texorpdfstring{Becker on Socialization: \emph{Becoming A
Marihuana
User}}{Becker on Socialization: Becoming A Marihuana User}}\label{becker-on-socialization-becoming-a-marihuana-user}

To an extent, both of these theories are synthesized and worked through
by Howard Becker, in his famous essay ``Becoming a Marihuana User.'' The
essay's starting point is to push back against the idea that Marijuana
is \emph{inherently} addictive, and instead to emphasize that ``getting
high'' is the outcome of a process of socialization. Specifically,
people must learn

\begin{enumerate}
\def\labelenumi{\arabic{enumi}.}
\tightlist
\item
  How ``to smoke the drug in a way that will produce real effects,''
\item
  How ``to recognize the effects and connect them with drug use,'' and
\item
  How ``to enjoy the sensations {[}you{]} perceive.''
\end{enumerate}

Collectively, this emphasizes that ``being high'' is a \emph{social}
category that requires socialization and not an inherent
\emph{biological} (or neurobiological) category.

\subsubsection{\texorpdfstring{Examples of Becker: Getting Drunk and
High \emph{For
Science}}{Examples of Becker: Getting Drunk and High For Science}}\label{examples-of-becker-getting-drunk-and-high-for-science}

There are two direct examples that illustrate Becker's theory:
\href{https://www.youtube.com/watch?v=IRBAZJ4lF0U}{grandmas smoking weed
for the first time}, and young people
\href{https://www.youtube.com/watch?v=L5J62dfPq88}{trying moonshine for
the first time}.

\section{Interaction}\label{interaction}

As implied by the previous discussion, people are socialized
fundamentally through \textbf{interaction}, or the exchange of meaning
with one another. Think about it this way: given the amount of your
lives that are organized through meaning, and that the vehicle of
meaning is symbols and language, and \emph{also} that the links between
what is doing the representation (the signifier) and what is being
represented (the signified) are (even at their most direct) ambiguous,
and (and the least direct) totally arbitrary, how else could you
understand the meanings of symbols except for repeatedly interacting
with others and the world as a whole?

\marginnote{\begin{footnotesize}

\textbf{Interaction} is the process through while people exchange
symbols and meaning with one another and the world in given social
settings.

\end{footnotesize}}

\subsection{The Guts of Interaction: Setting, Self-Identity,
Social-Identity}\label{the-guts-of-interaction-setting-self-identity-social-identity}

Any interaction is composed of three things. First, it must take place
in a \textbf{social setting}---that is, at a given time and place. The
setting carries with it assumptions about what norms apply, and what
social roles you and others can be expected to play. Within that
setting, an interaction is also composed of social and self-identities
giving off cues to one another in that interaction. Consequently, what
sometimes seem to be mundane activities in fact contain intricately
complex social interactions.

\subsection{``Normal'' Interaction}\label{normal-interaction}

The vast majority of our interactins are ``normal,'' which is to say
that, once we are adequately socialized, they take on the quality of
background aspects of our social lives.

In these ``normal'' interactions, we attend to people's facial
expressions, gestures, and eye contact, in ways that are shaped by our
own understanding of the nature of the interaction. In some
interactions---riding an elevator together with strangers, for
instance---we practice \textbf{civil inattention}, which is to say we
are aware of the presence of others but studiously avoid a direct,
attentive interaction with them. (In other words, while we know the
strangers are there in the elevator, we put our headphones on, pull out
our phones, and avoid making lingering eye contact or initating a direct
conversation.) Likewise, in normal interactions, we monitor the personal
space we have with others

\subsection{``Eruptions'' of
Interaction}\label{eruptions-of-interaction}

\subsection{Unequal Interactions}\label{unequal-interactions}

\subsection{Mediated Interaction}\label{mediated-interaction}

Consider this moment right now, in which I am writing these words and
you are reading them. Strictly speaking, this is one social moment, but
it is two moments in time. It is two moments in time in the sense that I
am typing these words at one time, and you are reading it at another,
but from the point of view of our symbolic exchange---enabled by
writing---those moments in time are collapsed together, such that my
position in writing these words and you in reading them, in the context
of the class, are the relevant difference!

The fact that this fantastic social moment is happening across time and
space is facilitated by the \textbf{mediation} of our interaction. That
is, there is a technological apparatus sitting between the direct,
face-to-face cues that we can pay attention to, cue on as part of the
taken-for-granted background of our everyday, well-socialized lives, or,
in the case of civil inattention, politely ignore.




\end{document}
